% ============================================================
%  04 · Signalschätzung / Estimation
%      Quelle: 04_signal_estimation.pdf + 01_Grundlagen.pdf
% ============================================================

\bereich[cSchaetz]{Schätztheorie — Grundbegriffe}
\begin{formelreihe}
  \infoblock{Schätzer: Algorithmus, der aus Beobachtungsvektor $\mathbf{X}$
    (Zufallsvariable) einen Schätzwert $\hat{\theta}$ für den
    deterministischen Parametervektor $\boldsymbol{\theta}$ liefert.} &
  \formelblock{Schätzfehler}{e=\theta-\hat{\theta}} &
  \formelblock{MSE (Skalar)}{\mathrm{MSE}=E\{e^2\}=E\{(\theta-\hat{\theta})^2\}} \\[0.3em]
  \formelblock{MSE (Vektor)}{\mathrm{MSE}=E\{\mathbf{e}\cdot\mathbf{e}^T\}
    =E\{(\boldsymbol{\theta}-\hat{\boldsymbol{\theta}})
    (\boldsymbol{\theta}-\hat{\boldsymbol{\theta}})^T\}} &
  \formelblock{Entwurfsziel}{\mathrm{MSE}\to\min\quad\text{(MMSE)}} &
  \formelblock{MSE-Zerlegung}{\mathrm{MSE}=\mathrm{Bias}^2+\mathrm{Var}(\hat{\theta})} \\
\end{formelreihe}
\bereichende

\bereich[cSchaetz]{Bewertungskriterien \& Eigenschaften}
\begin{formelreihe}
  \formelblock{Bias (syst.\ Fehler)}{b=E\{\hat{\theta}\}-\theta} &
  \formelblock{Kovarianz der Schätzung}{C_{\hat{\theta}}
    =E\!\left\{(\hat{\theta}{-}E\{\hat{\theta}\})
    (\hat{\theta}{-}E\{\hat{\theta}\})^T\right\}} &
  \formelblock{Erwartungstreu (unbiased)}{E\{\hat{\theta}\}=\theta
    \;\Leftrightarrow\; b=0} \\[0.3em]
  \formelblock{Asympt.\ erwartungstreu}{\lim_{N\to\infty}E\{\hat{\theta}\}=\theta} &
  \formelblock{Konsistent}{\lim_{N\to\infty}P(|\hat{\theta}-\theta|>\epsilon)=0} &
  \formelblock{Effizient}{C_{\hat{\theta}^*}\ge C_{\hat{\theta}_{\min}}\;\forall\,\hat{\theta}^*} \\
\end{formelreihe}
\bereichende

\bereich[cSchaetz]{Mittelwert- \& Varianzschätzer}
\begin{formelreihe}
  \formelblock{Stichprobenmittel (erw.-treu)}{\hat{\mu}_X
    =\dfrac{1}{N}\sum_{i=1}^{N}x_i,\quad E\{\hat{\mu}_X\}=\mu_X} &
  \formelblock{Varianz — biased}{\hat{\sigma}_{X,*}^2
    =\dfrac{1}{N}\sum_{i=1}^{N}(x_i-\hat{\mu}_X)^2} &
  \beispielblock{Bias der biased Varianz}{E\{\hat{\sigma}_{X,*}^2\}
    =\sigma_X^2\cdot\dfrac{N-1}{N}\;\neq\sigma_X^2} \\[0.3em]
  \formelblock{Varianz — unbiased}{\hat{\sigma}_X^2
    =\dfrac{1}{N-1}\sum_{i=1}^{N}(x_i-\hat{\mu}_X)^2} &
  \formelblock{Korrektur}{\hat{\sigma}_X^2
    =\dfrac{N}{N-1}\,\hat{\sigma}_{X,*}^2,\quad
    E\{\hat{\sigma}_X^2\}=\sigma_X^2} &
  \infoblock{Biased: asympt.\ erw.-treu\\
    ($\lim_{N\to\infty}\tfrac{N-1}{N}=1$).\\
    $\mu_X$ bekannt $\Rightarrow$ Nenner $N$ korrekt.} \\
\end{formelreihe}
\bereichende

\bereich[cSchaetz]{Korrelationsschätzer}
\begin{formelreihe}
  \formelblock{Kreuzkorrelation (def.)}{\varphi_{xy}[\kappa]
    =E\{X[k{+}\kappa]\cdot Y[k]\}} &
  \formelblock{Unnormierte Summe}{\hat{c}_{xy}[\kappa]
    =\sum_{k=0}^{N-\kappa-1}X[k{+}\kappa]\cdot Y[k]\quad(\kappa\ge0)} &
  \formelblock{Biased-Schätzer}{\hat{\varphi}_{xy}^*[\kappa]
    =\dfrac{1}{N}\,\hat{c}_{xy}[\kappa]} \\[0.3em]
  \formelblock{Unbiased-Schätzer}{\hat{\varphi}_{xy}[\kappa]
    =\dfrac{1}{N-|\kappa|}\,\hat{c}_{xy}[\kappa]} &
  \infoblock{Bevorzugt für kleine $|\kappa|$.\\
    Schätzfehler steigt für $|\kappa|\to N$.} &
  \\
\end{formelreihe}
\bereichende

\bereich[cSchaetz]{Wiener-Filter — Grundprinzip}
\begin{formelreihe}
  \infoblock{Beobachtetes Signal $x(t)$: Originalsignal $s(t)$\\
    verzerrt durch LTI $+$ Rauschen.\\[0.1em]
    Rekonstruktionsfilter $H$:\\
    minimiert Fehlerleistung $E\{|e|^2\}$.\\[0.1em]
    $e(t)=y(t)-s(t)$} &
  \formelblock{LDS des Fehlers}{\Phi_{ee}
    =\Phi_{xx}|H|^2-\Phi_{xs}H-\Phi_{sx}H^*+\Phi_{ss}} &
  \formelblock{Wiener-Filter (kontinuierlich)}{H(j\omega)
    =\dfrac{\Phi_{sx}(\omega)}{\Phi_{xx}(\omega)}
    =\dfrac{\Phi_{xs}^*(\omega)}{\Phi_{xx}(\omega)}} \\[0.3em]
  \formelblock{Wiener-Filter (zeitdiskret)}{H(e^{j\Omega})
    =\dfrac{\Phi_{sx}(\Omega)}{\Phi_{xx}(\Omega)}
    =\dfrac{\Phi_{xs}^*(\Omega)}{\Phi_{xx}(\Omega)}} &
  \beispielblock{Gestörter Kanal: $x=G*s+n$, $n\perp s$}{%
    \Phi_{xx}=\Phi_{ss}|G|^2+\Phi_{nn}\\[0.15em]
    \Phi_{sx}=\Phi_{ss}\cdot G^*(j\omega)\\[0.15em]
    H=\dfrac{\Phi_{ss}\cdot G^*}{\Phi_{ss}|G|^2+\Phi_{nn}}} &
  \infoblock{Einschränkungen:\\
    -- unendl.\ Beobachtungsdauer nötig\\
    -- Prozesse müssen stationär sein\\
    -- Einschwingvorgang\\[0.1em]
    $\to$ Kalman-Filter für endliche\\
    Beobachtung, zeitvariante Systeme} \\
\end{formelreihe}
\bereichende

\bereich[cSchaetz]{Kalman-Filter — Systemmodell}
\begin{formelreihe}
  \formelblock{Zustandsdifferenzengleichung}{x[k{+}1]=A[k]x[k]+B[k]u[k]+G[k]z[k]\\[0.1em]
    y[k]=C[k]x[k]+n[k]} &
  \infoblock{$x[k]$: Zustandsvektor (unbekannt)\\
    $u[k]$: bekannter Eingang\\
    $z[k]$: Systemrauschen (weiß, mw-frei)\\
    $n[k]$: Messrauschen (weiß, mw-frei)\\[0.1em]
    $\varphi_{zz}[\kappa,k]=\delta[\kappa]\cdot Z_0[k]$\\
    $\varphi_{nn}[\kappa,k]=\delta[\kappa]\cdot N_0[k]$} &
  \formelblock{Schätzfehler \& Ziel}{e[k]=\hat{x}[k]-x[k]\\[0.1em]
    F[k]=E\{e[k]\cdot e^T[k]\}\quad\text{(Fehlerkorr.-Matrix)}\\[0.1em]
    \mathrm{trace}\{F[k]\}\to\min} \\
\end{formelreihe}
\bereichende

\bereich[cSchaetz]{Kalman-Filter — Algorithmus (ohne Innovationsterm)}
\begin{formelreihe}
  \formelblock{Kalman-Gain \& Inno.-Kovarianz}{R[k]=N_0[k]+CF[k]C^T\\[0.15em]
    K[k]=AF[k]C^T\cdot R^{-1}[k]} &
  \formelblock{Update Zustandsschätzung}{\hat{x}[k{+}1]=A\hat{x}[k]+Bu[k]
    +K[k](y[k]-\hat{y}[k])\\[0.1em]
    \hat{y}[k]=C\hat{x}[k]} &
  \formelblock{Update Fehlerkorr.-Matrix}{F[k{+}1]=AF[k]A^T\\
    \quad-K[k]R[k]K^T[k]+GZ_0G^T\\[0.1em]
    \text{Start: }F[0]=0} \\
\end{formelreihe}
\bereichende

\bereich[cSchaetz]{Kalman-Filter — mit Innovationsterm (a-posteriori)}
\begin{formelreihe}
  \formelblock{Gewichtungsmatrix \& Gain}{R[k]=N_0[k]+CF[k]C^T\\[0.15em]
    M[k]=F[k]C^T\cdot R^{-1}[k]\\[0.1em]
    K[k]=A\cdot M[k]} &
  \formelblock{A-posteriori Zustandsschätzung}{\hat{x}^+[k]=\hat{x}[k]+M[k](y[k]-\hat{y}[k])\\[0.15em]
    \hat{x}[k{+}1]=A\hat{x}^+[k]+Bu[k]} &
  \formelblock{Update Fehlerkorr.-Matrix}{F[k{+}1]=AF[k]A^T\\
    \quad-AM[k]R[k]M^T[k]A^T+GZ_0G^T\\[0.1em]
    \text{Start: }F[0]=0} \\[0.3em]
  \beispielblock{Redundante Messung: $y_i=x+n_i$, $\hat{x}=k_1y_1+k_2y_2$}{%
    k_1+k_2=1\quad\text{(Erw.-treue)}\\
    k_1=\dfrac{\sigma_{n_2}^2}{\sigma_{n_1}^2+\sigma_{n_2}^2},\quad
    k_2=\dfrac{\sigma_{n_1}^2}{\sigma_{n_1}^2+\sigma_{n_2}^2}} &
  \infoblock{Eigenschaften:\\
    -- Einschwingvorgang des Gains\\
    -- LTI + stationäres Rauschen\\
    $\Rightarrow$ stationärer Endwert $K_\infty$\\
    -- In diesem Fall: Kalman $=$ Wiener} &
  \formelblock{Stationärer Endwert $f_\infty$ (Skalar)}{%
    f_\infty=\dfrac{A^2 N_0\,f_\infty}{N_0+C^2 f_\infty}+G Z_0 G^T\\[0.15em]
    \text{quadr.\ Gl.\ lösen, }R_\infty{=}N_0{+}C^2 f_\infty\\[0.1em]
    \Rightarrow K_\infty=\dfrac{A\,C\,f_\infty}{R_\infty}} \\
\end{formelreihe}
\bereichende

% ==========================================================
%  Kalman / Zustandsraum — Rezepte und durchgerechnete Beispiele
% ==========================================================

\bereich[cSchaetz]{REZEPT C1 · Zustandsraum aus Blockschaltbild ablesen}
\noindent\footnotesize
Jeder Verzögerer $z^{-1}$ ist ein Zustand $x_i$: sein \emph{Ausgang} ist $x_i[k]$,
sein \emph{Eingang} ist $x_i[k{+}1]$. Stelle für jeden Verzögerer-Eingang die
Summe aller einlaufenden Pfade auf.
\begin{rezept}
  \rs{Verzögerer durchnummerieren $x_1\dots x_n$ (Ausgang $=x_i[k]$).}
  \rs{Für jeden Eingang $x_i[k{+}1]$: alle einlaufenden Zweige addieren
      (Rückkopplungsfaktor $\to$ Eintrag in $\mathbf A$, Pfad von $u[k]\to\mathbf B$,
      Pfad von $z[k]\to\mathbf G$).}
  \rs{Ausgangsgleichung $y[k]=\mathbf C\,\mathbf x[k]+n[k]$: welcher Zustand (mit
      welchem Faktor) liegt am Messpunkt? $\to\mathbf C$. Additives Rauschen $n[k]$
      am Ausgang.}
  \rs{Reihenfolge im Vektor frei, aber konsistent halten (Eingang von $x_i$ links,
      Ausgang von $x_i$ rechts vom $z^{-1}$).}
\end{rezept}
\fallstrick{Skalierungsblock \emph{nach} dem Verzögerer (z.\,B.\ $\cdot 2$ am
  Messabgriff) gehört in $\mathbf C$, nicht in $\mathbf A$ (Fall $C{\neq}1$).\;
  Rauschpfad $z[k]$ läuft i.\,d.\,R.\ über einen eigenen Faktor in $\mathbf G$ —
  nicht mit $\mathbf B$ verwechseln.}
\miniexample{System 1.\ Ordnung}{%
  $x[k{+}1]=0{,}8\,x[k]+0{,}2\,z[k]+0{,}7\,u[k]$,\quad $y[k]=x[k]+n[k]$\\
  $\Rightarrow A{=}0{,}8,\;B{=}0{,}7,\;G{=}0{,}2,\;C{=}1$.}
\bereichende

\bereich[cSchaetz]{REZEPT C2 · Kalman-Algorithmus iterativ rechnen}
\noindent\footnotesize
Beide Varianten kommen vor. Start stets $\mathbf F[0]=\mathbf 0$, $\hat{\mathbf x}[0]=\mathbf 0$
(falls nicht anders angegeben). Pro Schritt $k$ in dieser Reihenfolge:
\begin{formelreihe}
  \parbox[t]{\linewidth}{%
    {\scriptsize\scshape\color{gray!65!black} Variante A — OHNE Innovationsterm
      (a-priori $\hat x[k]$)}\par\vspace{0.1em}
    \begin{rezept}
      \rs{$R[k]=N_0+C\,F[k]\,C^{T}$}
      \rs{$K[k]=A\,F[k]\,C^{T}\,R^{-1}[k]$}
      \rs{$\hat x[k{+}1]=A\,\hat x[k]+B\,u[k]+K[k]\,(y[k]-C\,\hat x[k])$}
      \rs{$F[k{+}1]=A\,F[k]\,A^{T}-K[k]\,R[k]\,K^{T}[k]+G\,Z_0\,G^{T}$}
    \end{rezept}} &
  \parbox[t]{\linewidth}{%
    {\scriptsize\scshape\color{gray!65!black} Variante B — MIT Innovationsterm
      (a-posteriori $\hat x^{+}[k]$)}\par\vspace{0.1em}
    \begin{rezept}
      \rs{$R[k]=N_0+C\,F[k]\,C^{T}$}
      \rs{$M[k]=F[k]\,C^{T}\,R^{-1}[k]$\quad(Gewichtungsmatrix),\;\; $K[k]=A\,M[k]$}
      \rs{$\hat x^{+}[k]=\hat x[k]+M[k]\,(y[k]-C\,\hat x[k])$}
      \rs{$\hat x[k{+}1]=A\,\hat x^{+}[k]+B\,u[k]$}
      \rs{$F[k{+}1]=A\,F[k]\,A^{T}-K[k]\,R[k]\,K^{T}[k]+G\,Z_0\,G^{T}$}
    \end{rezept}} &
  \infoblock{\textbf{Welche Variante?}\\
    „Schätzwert $\hat x[k]$ / Eingang $y[k]$, Ausgang $\hat x$" ohne extra
    Korrekturpfad $\to$ \emph{ohne}.\\
    Innovationspfad $e[k]{=}y{-}\hat y$ wirkt direkt auf $\hat x^{+}[k]$
    $\to$ \emph{mit}.\\[0.3em]
    \textbf{Alternative Nomenklatur} möglich: $\Sigma\,{\hat=}\,M$,
    $\Lambda\,{\hat=}\,$Rückführung, $\Phi,\Psi\,{\hat=}\,A,$ Verzög.\ —
    Struktur identisch.\\[0.3em]
    Nur $K[k]$ bzw.\ $M[k]$ und $F[k]$ sind \emph{zeitvariant}; $A,B,C,G$ konstant.}
\end{formelreihe}
\fallstrick{Erst $R$, dann $K/M$, dann $\hat x$, zuletzt $F$ — Reihenfolge nicht
  vertauschen.\; $F[0]{=}0\Rightarrow K[0]{=}0\Rightarrow F[1]=G Z_0 G^{T}$ (immer!).\;
  Bei $C{\neq}1$: $C F C^{T}$ und $C\hat x$ nicht vergessen.\;
  $u[k]=\cos(k\pi)=(-1)^k$ erzeugt die Folge $1,-1,1,\dots$}
\bereichende

\bereich[cSchaetz]{REZEPT C3 · Dimensionen \& stationärer Endwert}
\begin{formelreihe}
  \parbox[t]{\linewidth}{%
    {\scriptsize\scshape\color{gray!65!black} Dimensionen ($n$ Zustände, $1$ Ausgang)}\\[0.1em]
    {\footnotesize
    $\mathbf A:n{\times}n$,\;\; $\mathbf C:1{\times}n$,\;\;
    $\mathbf G:n{\times}q$,\;\; $\mathbf F:n{\times}n$,\\[0.15em]
    $R:1{\times}1$,\;\; $\mathbf M[k]:n{\times}1$,\;\; $\mathbf K[k]:n{\times}1$.\\[0.1em]
    \itshape ($q$ = Anzahl Systemrausch-Quellen $z$.)}} &
  \parbox[t]{\linewidth}{%
    {\scriptsize\scshape\color{gray!65!black} Stationärer Endwert (Skalar)}\par\vspace{0.05em}
    \begin{rezept}
      \rs{Setze $F[k{+}1]=F[k]=f_\infty$ in die $F$-Rekursion ein.}
      \rs{$K/M$ durch $f_\infty$ ausdrücken und einsetzen.}
      \rs{Auf Form $f_\infty^2+p\,f_\infty+q=0$ bringen, quadr.\ Gl.\ lösen.}
      \rs{\emph{Nur positive} Lösung $f_\infty$ behalten $\to$ $K_\infty$ bzw.\ $M_\infty$.}
    \end{rezept}} &
  \infoblock{\textbf{Existiert ein Endwert?} Ja, wenn System LTI \emph{und}
    Rauschen (schwach) stationär ist — dann Begründung mit genau diesem Satz.\\[0.3em]
    Mit den Endwerten initialisiert: \emph{kein} Einschwingvorgang, Filter
    zeitinvariant $\Rightarrow$ entspricht dem Wiener-Filter.}
\end{formelreihe}
\miniexample{Skalar, ohne Innovation ($A{=}0{,}8$, $C{=}1$, $G{=}0{,}2$, $Z_0{=}3$, $N_0{=}2$)}{%
  $R[k]{=}2{+}f[k]$,\quad $K[k]{=}\dfrac{0{,}8\,f[k]}{2{+}f[k]}$,\quad
  $f[k{+}1]{=}0{,}64\,f[k]-\dfrac{0{,}64\,f^{2}[k]}{2{+}f[k]}+0{,}12$\\[0.25em]
  \begin{tabular}{@{}l|cccc@{}}
    $k$ & 0 & 1 & 2 & 3\\\hline
    $f[k]$ & 0 & 0{,}12 & 0{,}1925 & 0{,}2324\\
    $K[k]$ & 0 & 0{,}0453 & 0{,}0702 & 0{,}0833
  \end{tabular}\quad
  $\hat x[k{+}1]{=}0{,}8\hat x{+}0{,}7u{+}K(y{-}\hat x)$\\[0.25em]
  Endwert: $f_\infty{=}0{,}64f_\infty-\frac{0{,}64f_\infty^2}{2+f_\infty}+0{,}12$
  $\Rightarrow f_\infty^2-1{,}73f_\infty-4{,}9{=}0$, $f_\infty{=}3{,}24$,
  $K_\infty{=}\frac{0{,}8\cdot3{,}24}{2+3{,}24}{=}0{,}495$.}
\miniexample{Matrix, 3.\ Ordnung, ohne Innovation (erster Schritt)}{%
  $\mathbf A{=}\!\begin{psmallmatrix}0&0&0\\1&0{,}7&0\\0&0&0{,}81\end{psmallmatrix}$,\;
  $\mathbf G{=}\!\begin{psmallmatrix}0&0\\1&0\\0&0{,}9\end{psmallmatrix}$,\;
  $\mathbf C{=}\!\begin{psmallmatrix}0&2&2\end{psmallmatrix}$,\;
  $\mathbf Z_0{=}\!\begin{psmallmatrix}1&0\\0&2\end{psmallmatrix}$,\; $R_0{=}3$.\\[0.2em]
  $\mathbf F[0]{=}\mathbf 0\Rightarrow R[0]{=}R_0{=}3\Rightarrow \mathbf K[0]{=}\mathbf 0$
  $\Rightarrow \mathbf F[1]{=}\mathbf G\,\mathbf Z_0\,\mathbf G^{T}
  {=}\!\begin{psmallmatrix}0&0&0\\0&1&0\\0&0&1{,}62\end{psmallmatrix}$.\quad
  $\dim\mathbf K[1]=3{\times}1$.}
\bereichende

\bereich[cSchaetz]{REZEPT C4 · Wiener-Filter als MSE-Minimierung im Zeitbereich}
\noindent\footnotesize
Tritt auf, wenn das Filter eine \emph{Konstante} $h_0$ ist (kein LDS gefordn.).
Ziel: $E\{e^2\}\to\min$.
\begin{rezept}
  \rs{Schätzwert hinschreiben: $\hat s(t)=h_0\cdot x(t)$, dabei $x(t)$ durch
      $s$, verzögerte $s$ und Rauschen ausdrücken.}
  \rs{Fehler $e(t)=s(t)-\hat s(t)$ einsetzen und nach Termen ordnen.}
  \rs{$E\{e^2\}$ ausmultiplizieren; nutze $E\{s(t)s(t{-}\tau)\}=\varphi_{ss}(\tau)$,
      $E\{s\cdot n\}=0$, $E\{n^2\}=\varphi_{nn}(0)$.}
  \rs{Ergebnis ist Parabel $E\{e^2\}=\alpha\,h_0^2+\beta\,h_0+\gamma$.}
  \rs{$\dfrac{\mathrm d}{\mathrm d h_0}E\{e^2\}=0 \Rightarrow h_0=-\dfrac{\beta}{2\alpha}$.}
\end{rezept}
\miniexample{$\hat s=h_0(5s(t)+n_1(t)+s(t{-}2))$, $\varphi_{ss}(0){=}0{,}6$, $\varphi_{nn}{=}10\delta$}{%
  $e=s(1{-}5h_0)-h_0 s(t{-}2)-h_0 n_1$\;
  $\Rightarrow E\{e^2\}=28{,}6\,h_0^2-6{,}6\,h_0+0{,}6$\;
  $\Rightarrow h_0=\dfrac{6{,}6}{2\cdot28{,}6}=0{,}115$.}
\bereichende
